\documentclass{BeamerTemplate}

\title{Chapitre 3 -- Estimations et tests d'hypothèse (1\textsuperscript{re} partie)}
\subtitle{STT-1920 Méthodes statistiques}
\date{}

% Styles des graphiques
\pgfplotsset{
  base/.style={
    tick label style={font=\scriptsize}, label style={font=\small},
    title style={font=\small\bfseries},
    yticklabel style={/pgf/number format/fixed, /pgf/number format/precision=2},
    scaled y ticks=false,
  }
}

% Commande pour dessiner une cible TikZ (qualité d'estimation)
\newcommand{\cibleTikz}[5]{% #1=titre, #2=dx, #3=dy, #4=dispersion, #5=label
  \begin{tikzpicture}[scale=0.32]
    \draw[thick, fill=gray!10] (0,0) circle (2);
    \draw[thick, fill=gray!20] (0,0) circle (1.4);
    \draw[thick, fill=Dominante!30] (0,0) circle (0.8);
    \draw[thick, fill=Accent!40] (0,0) circle (0.3);
    \draw[dashed, thin, gray] (-2.1,0) -- (2.1,0);
    \draw[dashed, thin, gray] (0,-2.1) -- (0,2.1);
    % Points d'impact
    \foreach \ang/\r in {30/0.2, 120/0.3, 210/0.15, 300/0.25, 75/0.18, 180/0.22, 260/0.28} {
      \fill[rougelaval] ({#2 + (\r*#4)*cos(\ang)}, {#3 + (\r*#4)*sin(\ang)}) circle (2pt);
    }
    \node[below, font=\tiny\bfseries] at (0,-2.1) {#1};
  \end{tikzpicture}%
}

\begin{document}

\begin{frame}[t]
  \titlepage
\end{frame}

%===============================================================================
\section{Introduction à l'inférence statistique}
%===============================================================================

\begin{frame}[t]{Où en sommes-nous?}
  \begin{itemize}\itemsep 4mm
    \item Au \alert{chapitre 1}, nous avons appris à \textbf{décrire} un ensemble de données (moyenne, écart-type, médiane, histogrammes\ldots).
    \item Au \alert{chapitre 2}, nous avons étudié la \textbf{théorie des probabilités} en supposant les lois et leurs paramètres \emph{parfaitement connus}.
    \begin{itemize}
      \item \emph{Exemple :} la masse des oranges suit une loi $\mathcal{N}(\mu = 204,\, \sigma^2 = 256)$.
    \end{itemize}
  \end{itemize}

  \bigskip
  \uncover<2->{
  \begin{Boite}{Le problème en pratique}
    Dans la vraie vie, on ne connaît \alert{jamais} la vraie distribution ni les vraies valeurs des paramètres $\mu$, $\sigma^2$ ou $p$ de la population !
  \end{Boite}
  }

  \medskip
  \uncover<3->{
    $\Rightarrow$ C'est l'objet de l'\alert{inférence statistique} : tirer des conclusions sur une \textbf{population} à partir d'un \textbf{échantillon}.
  }
\end{frame}

\begin{frame}[t]{Démarche générale}
  \begin{enumerate}\itemsep 6mm
    \item<1-> On recueille un \alert{échantillon aléatoire} de taille $n$ dans la population d'intérêt.
    \item<2-> On postule (selon la théorie ou l'analyse descriptive) un \alert{modèle probabiliste} adapté au phénomène (loi normale, binomiale, Poisson\ldots).
    \item<3-> Ce modèle comporte des \alert{paramètres inconnus} ($\mu$, $\sigma^2$, $p$\ldots) : on doit les \textbf{estimer} à partir des données observées.
  \end{enumerate}

  \bigskip
  \uncover<4->{
  \begin{Boite}{Deux grandes familles d'inférence}
    \begin{itemize}
      \item \textbf{L'estimation} (ponctuelle et par intervalle de confiance)
      \item \textbf{Les tests d'hypothèses} (prise de décision face à une affirmation)
    \end{itemize}
  \end{Boite}
  }
\end{frame}

\begin{frame}[t]{Exemples motivants}
  \begin{itemize}\itemsep 6mm
    \item \textbf{Biologie cellulaire :} Le diamètre des cellules hépatiques suit une loi $\mathcal{N}(\mu, \sigma^2)$.
      \begin{itemize}
        \item On mesure 20 cellules : comment estimer le diamètre moyen $\mu$ et la variabilité $\sigma^2$ de toutes les cellules de ce type?
      \end{itemize}

    \item \textbf{Épidémiologie animale :} Le nombre de veaux infectés parmi $n$ veaux examinés suit une loi $\text{Binomiale}(n, p)$.
      \begin{itemize}
        \item Sur 80 veaux examinés, 12 présentent des symptômes. Quelle est la proportion réelle $p$ d'infection dans le troupeau?
      \end{itemize}

    \item \textbf{Agronomie :} Deux engrais sont testés sur des parcelles. L'engrais A donne-t-il un rendement moyen supérieur à l'engrais B?
  \end{itemize}
\end{frame}

%===============================================================================
\section{Vocabulaire et définitions}
%===============================================================================

\begin{frame}[t]{Population vs Échantillon}
  \begin{columns}[T]
    \column{0.48\textwidth}
      \begin{Boite}{Population (inaccessible)}
        \begin{itemize}
          \item Ensemble complet de toutes les unités d'intérêt.
          \item Caractérisée par des \alert{paramètres} fixes mais inconnus.
          \item Notations : $\mu$, $\sigma^2$, $\sigma$, $p$.
        \end{itemize}
      \end{Boite}

    \column{0.48\textwidth}
      \begin{Boite}{Échantillon (observé)}
        \begin{itemize}
          \item Sous-ensemble de $n$ unités tirées au hasard.
          \item Fournit des \alert{statistiques} calculées sur les données.
          \item Notations : $\overline{x}$, $s^2$, $s$, $\hat{p}$.
        \end{itemize}
      \end{Boite}
  \end{columns}

  \medskip
  \begin{center}
  \begin{tikzpicture}[scale=0.8]
    \node[draw=Dominante, fill=Dominante!15, rounded corners, inner sep=5pt, text width=4.3cm, align=center] (pop) at (0,0)
      {\textbf{Population}\\Paramètres : $\mu, \sigma^2, p$ (\emph{inconnus})};
    \node[draw=Accent, fill=Accent!15, rounded corners, inner sep=5pt, text width=4.3cm, align=center] (ech) at (6.2,0)
      {\textbf{Échantillon ($n$)}\\Statistiques : $\bar{x}, s^2, \hat{p}$ (\emph{calculées})};

    \draw[->, thick, Dominante] (pop.north) to[bend left=25] node[above, font=\scriptsize\bfseries, color=Dominante] {Échantillonnage aléatoire} (ech.north);
    \draw[->, thick, Accent] (ech.south) to[bend left=25] node[below, font=\scriptsize\bfseries, color=Accent] {Inférence statistique} (pop.south);
  \end{tikzpicture}
  \end{center}
\end{frame}

\begin{frame}[t]{Variables aléatoires vs Observations}
  \begin{Boite}{Distinction entre majuscules et minuscules}
    \begin{itemize}\itemsep 2mm
      \item \textbf{Avant l'expérience :} les mesures $X_1, X_2, \ldots, X_n$ sont des \alert{variables aléatoires}.
      \item \textbf{Après l'expérience :} les valeurs obtenues $x_1, x_2, \ldots, x_n$ sont des \alert{nombres réels} fixés.
    \end{itemize}
  \end{Boite}

  \medskip
  \uncover<2->{
  \begin{Boite}{Estimateur vs Estimation ponctuelle}
    \begin{itemize}\itemsep 2mm
      \item \textbf{Estimateur} $T = g(X_1, \ldots, X_n)$ : \alert{variable aléatoire} (formule générale).
      \item \textbf{Estimation} $t = g(x_1, \ldots, x_n)$ : \alert{valeur numérique} sur l'échantillon observé.
    \end{itemize}
  \end{Boite}
  }

  \medskip
  \uncover<3->{
    $\Rightarrow$ Un estimateur a donc une \textbf{loi de probabilité}, une \textbf{espérance} et une \textbf{variance} !
  }
\end{frame}

\begin{frame}[t]{Tableau synthèse des notations}
  \centering
  \renewcommand{\arraystretch}{1.3}
  \begin{tabular}{lcccc}\toprule
    \textbf{Paramètre} & \textbf{Estimateur} & \textbf{Estimation} & \textbf{Espérance} & \textbf{Erreur-type} \\
    (population) & (v.a.) & (valeur) & $E[T]$ & $\sigma_T = \sqrt{\text{Var}(T)}$ \\ \midrule
    Moyenne $\mu$ & $\overline{X} = \frac{1}{n}\sum X_i$ & $\bar{x} = \frac{1}{n}\sum x_i$ & $\mu$ & $\frac{\sigma}{\sqrt{n}}$ \\[3mm]
    Variance $\sigma^2$ & $S^2 = \frac{1}{n-1}\sum (X_i - \overline{X})^2$ & $s^2$ & $\sigma^2$ & $\sqrt{\frac{2\sigma^4}{n-1}}$\footnotemark \\[3mm]
    Proportion $p$ & $\widehat{P} = \frac{Y}{n}$ & $\hat{p} = \frac{y}{n}$ & $p$ & $\sqrt{\frac{p(1-p)}{n}}$ \\ \bottomrule
  \end{tabular}
  \footnotetext{Sous l'hypothèse de normalité de la population.}

  \bigskip
  \uncover<2->{
  \begin{Boite}{Erreur-type}
    L'\alert{erreur-type} d'un estimateur est l'écart-type de sa distribution d'échantillonnage : $\text{ET}(T) = \sqrt{\text{Var}(T)}$. Elle mesure la \textbf{précision} de l'estimateur.
  \end{Boite}
  }
\end{frame}

%===============================================================================
\section{Propriétés d'un bon estimateur}
%===============================================================================

\begin{frame}[t]{Qualités d'un bon estimateur}
  On souhaite qu'un estimateur possède deux qualités fondamentales :
  \begin{columns}[T]
    \column{0.48\textwidth}
      \begin{Boite}{1. Être sans biais (juste)}
        \small
        En moyenne, il vise dans le mille :
        $$E[T] = \theta \quad (\text{Biais} = 0)$$
      \end{Boite}

    \column{0.48\textwidth}
      \begin{Boite}{2. Être précis (faible variance)}
        \small
        Les estimations varient très peu :
        $$\text{Var}(T) \quad \text{est minimale}$$
      \end{Boite}
  \end{columns}

  \smallskip
  \uncover<2->{
  \centering
  \begin{tabular}{cc}
    \cibleTikz{Sans biais et précis}{0}{0}{0.9}{} &
    \cibleTikz{Biaisé mais précis}{0.9}{0.7}{0.9}{} \\
    \cibleTikz{Sans biais mais imprécis}{0}{0}{2.4}{} &
    \cibleTikz{Biaisé et imprécis}{0.9}{0.7}{2.4}{}
  \end{tabular}
  }
\end{frame}

\begin{frame}[t]{Le biais : attention aux confusions !}
  Il faut distinguer plusieurs sens du mot « biais » :
  \begin{itemize}\itemsep 4mm
    \item \textbf{Biais de sélection / d'échantillonnage :}
      \begin{itemize}
        \item L'échantillon n'est pas représentatif de la population cible (ex. : sondage téléphonique auprès des abonnés à un journal spécifique).
        \item C'est une erreur de \alert{protocole} ou de collecte, impossible à corriger mathématiquement par la suite !
      \end{itemize}

    \item \textbf{Biais mathématique de l'estimateur :}
      \begin{itemize}
        \item $\text{Biais}(T) = E[T] - \theta$.
        \item C'est une propriété de la formule mathématique elle-même.
        \item \emph{Exemple :} $\frac{1}{n}\sum (X_i - \overline{X})^2$ est biaisé (sous-estime $\sigma^2$ en moyenne). On corrige ce biais en divisant par $n-1$.
      \end{itemize}
  \end{itemize}
\end{frame}

\begin{frame}[t]{Comment agir sur la précision?}
  L'erreur-type mesure l'imprécision d'une estimation. Pour la réduire, on dispose de deux leviers principaux :

  \bigskip
  \begin{enumerate}\itemsep 5mm
    \item \textbf{Choisir un estimateur de faible variance :}
      \begin{itemize}
        \item Entre deux estimateurs sans biais, on préfère toujours celui qui a la plus petite variance (estimateur le plus efficace).
      \end{itemize}

    \item \textbf{Augmenter la taille de l'échantillon $n$ :}
      \begin{itemize}
        \item Pour la moyenne, $\text{ET}(\overline{X}) = \frac{\sigma}{\sqrt{n}}$.
        \item Pour \alert{doubler la précision} (diviser l'erreur-type par 2), il faut \alert{quadrupler} la taille de l'échantillon ($4n$) !
      \end{itemize}
  \end{enumerate}
\end{frame}

%===============================================================================
\section{Estimation de la moyenne \texorpdfstring{$\mu$}{mu}}
%===============================================================================

\begin{frame}[t]{Exemple 1 : Pression artérielle systolique}
  On suppose que la pression artérielle systolique (en mm\,Hg) d'une population d'adultes sains suit une loi normale $\mathcal{N}(\mu, \sigma^2)$.

  \medskip
  On mesure la pression chez $n = 7$ individus choisis au hasard :
  \begin{center}
    \textbf{141 \quad 130 \quad 135 \quad 138 \quad 137 \quad 122 \quad 135}
  \end{center}

  \bigskip
  \uncover<2->{
  \begin{Boite}{Estimation ponctuelle de la moyenne}
    $$\bar{x} = \frac{1}{n}\sum_{i=1}^n x_i = \frac{141 + 130 + 135 + 138 + 137 + 122 + 135}{7} = \frac{938}{7} = \mathbf{134\text{ mm\,Hg}}$$
  \end{Boite}
  }

  \bigskip
  \uncover<3->{
    \textbf{Question :} Si on avait choisi 7 autres individus, aurait-on obtenu exactement 134?
    \alert{Non !} La moyenne observée varie d'un échantillon à l'autre.
  }
\end{frame}

\begin{frame}[t]{Variabilité d'échantillonnage : une simulation}
  Supposons que la vraie population suit $\mathcal{N}(\mu = 133,\, \sigma^2 = 62)$. Si on tirait 10 échantillons indépendants de taille $n=7$ :

  \medskip
  \centering
  \small
  \begin{tabular}{ccccccccc}\toprule
    \textbf{Éch.} & $x_1$ & $x_2$ & $x_3$ & $x_4$ & $x_5$ & $x_6$ & $x_7$ & $\bar{x}$ \\ \midrule
    1 & 141 & 130 & 135 & 138 & 137 & 122 & 135 & \textbf{134{,}0} \\
    2 & 132 & 128 & 139 & 133 & 132 & 127 & 138 & \textbf{132{,}7} \\
    3 & 140 & 125 & 122 & 126 & 136 & 128 & 131 & \textbf{129{,}7} \\
    4 & 131 & 137 & 136 & 130 & 140 & 134 & 124 & \textbf{133{,}1} \\
    5 & 143 & 125 & 131 & 121 & 139 & 150 & 131 & \textbf{134{,}3} \\
    6 & 132 & 140 & 128 & 136 & 141 & 135 & 140 & \textbf{136{,}0} \\
    7 & 135 & 143 & 134 & 133 & 138 & 123 & 134 & \textbf{134{,}3} \\ \bottomrule
  \end{tabular}

  \bigskip
  \uncover<2->{
    $\Rightarrow$ Les moyennes s'éparpillent autour de la vraie moyenne ($\mu = 133$).
  }
\end{frame}

\begin{frame}[t]{Distribution de la population vs distribution de $\overline{X}$}
  \begin{center}
  \begin{tikzpicture}
    \begin{axis}[
      base, width=10.5cm, height=5.5cm,
      xmin=105, xmax=160, ymin=0, ymax=0.16,
      xlabel={Pression (mm\,Hg)}, ylabel={Densité},
      legend pos=north east, legend style={font=\scriptsize}
    ]
      % Population: N(133, 62) => sigma = 7.87
      \addplot[thick, Dominante, domain=105:160, samples=100]
        {exp(-((x-133)^2)/(2*62)) / (sqrt(2*pi*62))};
      \addlegendentry{Population $X \sim \mathcal{N}(133, 62)$}

      % Moyenne échantillon: N(133, 62/7) => sigma_bar = sqrt(62/7) = 2.97
      \addplot[thick, Accent, domain=115:150, samples=120]
        {exp(-((x-133)^2)/(2*(62/7))) / (sqrt(2*pi*(62/7)))};
      \addlegendentry{Moyenne $\overline{X} \sim \mathcal{N}(133,\, 62/7)$}

      \draw[dashed, gray] (axis cs:133,0) -- (axis cs:133,0.14);
      \node[above, font=\scriptsize] at (axis cs:133,0.14) {$\mu = 133$};
    \end{axis}
  \end{tikzpicture}
  \end{center}

  \vspace{-2mm}
  \alert{Constat majeur :} $\overline{X}$ a la même moyenne $\mu$, mais une variance $n$ fois plus petite !
\end{frame}

\begin{frame}[t]{Propriétés théoriques de $\overline{X}$}
  Soit un échantillon aléatoire $X_1, \ldots, X_n$ i.i.d. d'espérance $\mu$ et de variance $\sigma^2$.

  \bigskip
  \uncover<1->{
  \begin{Boite}{1. Espérance (sans biais)}
    $$E[\overline{X}] = E\left[\frac{1}{n}\sum_{i=1}^n X_i\right] = \frac{1}{n}\sum_{i=1}^n E[X_i] = \frac{1}{n}(n\mu) = \alert{\mu}$$
  \end{Boite}
  }

  \bigskip
  \uncover<2->{
  \begin{Boite}{2. Variance et erreur-type}
    $$\text{Var}(\overline{X}) = \text{Var}\left(\frac{1}{n}\sum_{i=1}^n X_i\right) = \frac{1}{n^2}\sum_{i=1}^n \text{Var}(X_i) = \frac{n\sigma^2}{n^2} = \alert{\frac{\sigma^2}{n}}$$
    $$\text{ET}(\overline{X}) = \alert{\frac{\sigma}{\sqrt{n}}}$$
  \end{Boite}
  }
\end{frame}

%===============================================================================
\section{Estimation de la variance \texorpdfstring{$\sigma^2$}{sigma2}}
%===============================================================================

\begin{frame}[t]{Estimateur de la variance $\sigma^2$}
  On définit la variance échantillonnale par :
  \begin{Boite}{Variance échantillonnale $S^2$}
    $$S^2 = \frac{1}{n-1}\sum_{i=1}^n (X_i - \overline{X})^2$$
  \end{Boite}

  \bigskip
  \uncover<2->{
  \textbf{Pourquoi diviser par $n-1$ plutôt que par $n$?}
  \begin{itemize}\itemsep 3mm
    \item Les écarts se mesurent par rapport à $\overline{X}$ (estimé) et non par rapport au vrai $\mu$.
    \item Puisque $\sum (X_i - \overline{X}) = 0$, dès qu'on connaît $n-1$ écarts, le dernier est imposé : il y a \alert{$n-1$ degrés de liberté}.
    \item En divisant par $n-1$, l'estimateur devient \alert{sans biais} :
      $$E[S^2] = \sigma^2$$
  \end{itemize}
  }
\end{frame}

\begin{frame}[t]{Variance de $S^2$}
  Puisque $S^2$ est calculé à partir de l'échantillon, c'est aussi une variable aléatoire qui fluctue.

  \bigskip
  \begin{Boite}{Propriétés de $S^2$ (sous l'hypothèse de normalité)}
    \begin{itemize}\itemsep 3mm
      \item \textbf{Espérance :} $E[S^2] = \sigma^2$ (sans biais)
      \item \textbf{Variance :} $\text{Var}(S^2) = \frac{2\sigma^4}{n-1}$
      \item \textbf{Erreur-type :} $\text{ET}(S^2) = \sqrt{\frac{2\sigma^4}{n-1}} = \frac{\sqrt{2}\sigma^2}{\sqrt{n-1}}$
    \end{itemize}
  \end{Boite}

  \bigskip
  \uncover<2->{
    $\Rightarrow$ Plus $n$ est grand, plus $S^2$ se rapproche précisément de la vraie variance $\sigma^2$.
  }
\end{frame}

\begin{frame}[t]{Retour sur l'exemple 1 : Calculs complets}
  Échantillon de pression systolique ($n=7$) :
  \begin{center}
    141 \quad 130 \quad 135 \quad 138 \quad 137 \quad 122 \quad 135
  \end{center}

  \medskip
  \begin{enumerate}\itemsep 3mm
    \item<1-> \textbf{Moyenne échantillonnale :}
      $$\bar{x} = \frac{938}{7} = \mathbf{134\text{ mm\,Hg}}$$
    \item<2-> \textbf{Variance échantillonnale :}
      $$s^2 = \frac{1}{7-1} \left[(141-134)^2 + \ldots + (135-134)^2\right] = \frac{236}{6} = \mathbf{39{,}33\text{ (mm\,Hg)}^2}$$
    \item<3-> \textbf{Écart-type échantillonnal :}
      $$s = \sqrt{39{,}33} = \mathbf{6{,}27\text{ mm\,Hg}}$$
    \item<4-> \textbf{Erreur-type estimée de la moyenne :}
      $$\widehat{\text{ET}}(\overline{X}) = \frac{s}{\sqrt{n}} = \frac{6{,}27}{\sqrt{7}} = \mathbf{2{,}37\text{ mm\,Hg}}$$
  \end{enumerate}
\end{frame}

%===============================================================================
\section{Estimation d'une proportion \texorpdfstring{$p$}{p}}
%===============================================================================

\begin{frame}[t]{Exemple 2 : Un sondage d'opinion}
  Dans un sondage réalisé auprès de $n = 1028$ Québécois, on demande : \emph{« Quel enjeu influencera le plus votre vote? »}

  \medskip
  \textbf{Résultat :} $y = 216$ personnes répondent l'environnement.

  \bigskip
  \uncover<2->{
  \begin{Boite}{Modélisation}
    Soit $Y$ le nombre de répondants choisissant l'environnement parmi $n$.
    $$Y \sim \text{Binomiale}(n,\, p)$$
    où $p$ est la vraie proportion inconnue dans la population québécoise.
  \end{Boite}
  }

  \bigskip
  \uncover<3->{
    Comment estimer $p$ et quelle est la précision de notre résultat?
  }
\end{frame}

\begin{frame}[t]{Propriétés de l'estimateur de proportion}
  \begin{Boite}{Estimateur de la proportion}
    $$\widehat{P} = \frac{Y}{n}$$
  \end{Boite}

  \bigskip
  Puisque $Y \sim \text{Binomiale}(n, p)$, on sait que $E[Y] = np$ et $\text{Var}(Y) = np(1-p)$.

  \bigskip
  \uncover<2->{
  \begin{itemize}\itemsep 4mm
    \item \textbf{Espérance (sans biais) :}
      $$E[\widehat{P}] = E\left[\frac{Y}{n}\right] = \frac{1}{n} E[Y] = \frac{np}{n} = \alert{p}$$
    \item \textbf{Variance :}
      $$\text{Var}(\widehat{P}) = \text{Var}\left(\frac{Y}{n}\right) = \frac{1}{n^2}\text{Var}(Y) = \frac{np(1-p)}{n^2} = \alert{\frac{p(1-p)}{n}}$$
    \item \textbf{Erreur-type :} $\text{ET}(\widehat{P}) = \alert{\sqrt{\frac{p(1-p)}{n}}}$
  \end{itemize}
  }
\end{frame}

\begin{frame}[t]{Retour sur l'exemple 2 : Application}
  Données du sondage : $n = 1028$, $y = 216$.

  \bigskip
  \begin{itemize}\itemsep 5mm
    \item<1-> \textbf{Estimation ponctuelle de la proportion :}
      $$\hat{p} = \frac{216}{1028} \approx \mathbf{0{,}2101 \quad (21{,}0\%)}$$

    \item<2-> \textbf{Erreur-type estimée :}
      Puisque $p$ est inconnu, on remplace $p$ par $\hat{p}$ dans la formule :
      $$\widehat{\text{ET}}(\widehat{P}) = \sqrt{\frac{\hat{p}(1-\hat{p})}{n}} = \sqrt{\frac{0{,}2101 \times 0{,}7899}{1028}} \approx \mathbf{0{,}0127 \quad (1{,}27\%)}$$

    \item<3-> \textbf{Conclusion :} L'estimation est de $21{,}0\%$ avec une incertitude d'échantillonnage d'environ $1{,}3\%$.
  \end{itemize}
\end{frame}

\begin{frame}[t]{Conclusion de la partie 1}
  \begin{Boite}{Ce que nous savons faire à ce stade}
    \begin{itemize}\itemsep 3mm
      \item Calculer une \alert{estimation ponctuelle} pour une moyenne ($\bar{x}$), une variance ($s^2$) ou une proportion ($\hat{p}$).
      \item Évaluer sa précision grâce à son \alert{erreur-type}.
    \end{itemize}
  \end{Boite}

  \bigskip
  \uncover<2->{
  \begin{Boite}{Limite de l'estimation ponctuelle}
    Donner une valeur unique (ex. : $\bar{x} = 134$) n'indique pas à quel point nous sommes confiants que la vraie valeur $\mu$ se trouve à proximité.
  \end{Boite}
  }

  \bigskip
  \uncover<3->{
    $\Rightarrow$ \textbf{Prochaine étape (partie 2) :} construire des \alert{intervalles de confiance} pour quantifier formellement cette certitude.
  }
\end{frame}

\end{document}
